\chapter{Singularity Analysis}
\label{ch04:top}

Chapter~2 established the basic relationship between generating functions and asymptotics: if $A(z) = \sum_{n \ge 0} a_n z^n$ has radius of convergence $R$, then the sequence $a_n$ grows exponentially at rate $1/R$. That is the coarsest possible information about the coefficients. This chapter refines it considerably. The central insight, due to Flajolet and Odlyzko~\cite{FlajoletOdlyzko1990}, is that the \emph{type} of the dominant singularity of $A$ on its circle of convergence determines not only the exponential growth rate $\rho^{-n}$ (where $\rho = R$) but also the subexponential correction factor and the multiplicative constant. The full asymptotic takes the form $a_n \sim C n^{\alpha - 1} \rho^{-n}$, and reading off $C$, $\alpha$, and $\rho$ is a matter of inspecting the local behavior of $A$ near its singularity. This is the Flajolet--Odlyzko transfer theorem, the workhorse of analytic combinatorics~\cite{FlajoletSedgewick2009}.

\section{Exponential Growth Rate}

The Cauchy--Hadamard formula gives $1/R = \limsup_{n\to\infty} |a_n|^{1/n}$. Rephrasing: if $\rho = R$ is the radius of convergence, then $a_n$ grows no faster than $(\rho^{-1} + \varepsilon)^n$ for any $\varepsilon > 0$, and $\rho^{-n}$ is the sharpest exponential bound. More precisely, $|a_n| \le (\rho^{-1} + \varepsilon)^n$ eventually, and for any $C > 1$ the bound $|a_n| = \rho^{-n} \cdot o(C^n)$ holds. The factor $\rho^{-n}$ is called the \emph{exponential part} of the asymptotic. What $\rho$ does \emph{not} tell us is any polynomial or logarithmic correction: two sequences $n^{1/2} \rho^{-n}$ and $n^{-3/2} \rho^{-n}$ have the same exponential growth rate but vastly different behaviors for finite $n$. The polynomial correction is what singularity analysis supplies.

\section{Dominant Singularities}

\begin{definition}
A \emph{dominant singularity} of a power series $A(z)$ with radius of convergence $\rho$ is a point $\zeta$ with $|\zeta| = \rho$ at which $A$ fails to extend analytically.
\end{definition}

Every power series has at least one dominant singularity, since $A$ cannot be continued to an entire function unless it converges everywhere. The analytic behavior of $A$ near each dominant singularity contributes an oscillatory term to $a_n$. When there are multiple dominant singularities $\zeta_1, \ldots, \zeta_k$ on the circle $|z| = \rho$, the coefficient $a_n$ is a superposition of terms $c_j \zeta_j^{-n}$, and since $|\zeta_j^{-n}| = \rho^{-n}$ for all $j$, these terms interfere, producing oscillations rather than clean growth.

A simple illustration: the generating function $1/(1-z^2)$ has radius of convergence $1$, with two dominant singularities at $z = 1$ and $z = -1$. Partial fractions give
\[
\frac{1}{1-z^2} = \frac{1}{2}\left(\frac{1}{1-z} + \frac{1}{1+z}\right),
\]
so the coefficient of $z^n$ is $\tfrac{1}{2}(1 + (-1)^n)$, which is $1$ for even $n$ and $0$ for odd $n$. The coefficient sequence oscillates with period $2$ rather than growing monotonically, which is the characteristic signature of two dominant singularities with opposite signs. No single real exponential $C\rho^{-n}$ approximates these coefficients uniformly.

For the remainder of this chapter, we restrict to the \emph{aperiodic case}: $A$ has a unique dominant singularity, located at a positive real point $z = \rho$. This is the generic situation for combinatorial classes defined by positive structural rules, and it is where the transfer theorem applies cleanly.

\section{Rational Generating Functions}

When $A(z) = P(z)/Q(z)$ is rational with $\deg P < \deg Q$, the partial fraction decomposition gives an exact formula for every coefficient. If $Q$ has zeros $z_1, \ldots, z_r$ with multiplicities $m_1, \ldots, m_r$, then
\[
a_n = \sum_{i=1}^r \mathrm{Pol}_i(n) \cdot z_i^{-n},
\]
where $\mathrm{Pol}_i$ is a polynomial of degree $m_i - 1$ in $n$. The dominant contribution comes from the zero of $Q$ of smallest modulus, say $z_1 = \rho$ with multiplicity $m$. All other terms are $O(|z_i|^{-n}) = O((\rho + \delta)^{-n})$ for some $\delta > 0$, which is exponentially smaller. Consequently
\[
a_n \sim C n^{m-1} \rho^{-n}
\]
as $n \to \infty$, for an explicit constant $C$ determined by the Laurent expansion of $A$ at $z = \rho$.

\begin{example}[Fibonacci numbers]
Let $A(z) = z/(1 - z - z^2)$, the generating function for the Fibonacci sequence established in Chapter~1. The denominator $1 - z - z^2$ factors as $-(z - 1/\varphi)(z + \varphi)$, where $\varphi = (1+\sqrt{5})/2$ is the golden ratio. The two poles are at $z = 1/\varphi \approx 0.618$ and $z = -\varphi \approx -1.618$. Since $1/\varphi < \varphi$, the dominant pole is $z = 1/\varphi$, and it is simple, so $m = 1$.

The partial fraction decomposition was already carried out in Chapter~1 using the factorization $1 - z - z^2 = (1 - \varphi z)(1 - \hat\varphi z)$, where $\hat\varphi = (1 - \sqrt{5})/2$. The result is $A(z) = (1/\sqrt{5})/(1 - \varphi z) - (1/\sqrt{5})/(1 - \hat\varphi z)$, giving $C = 1/\sqrt{5}$. The transfer from the rational case then gives
\[
F_n \sim \frac{\varphi^n}{\sqrt{5}},
\]
which is Binet's formula to leading order. (The subleading term from the pole at $-\varphi$ is $(-1)^{n+1}/(\sqrt{5}\,\varphi^n)$, which decays exponentially and contributes less than $1/2$ in absolute value for all $n \ge 1$, so the round-to-nearest interpretation recovers the exact formula.)
\end{example}

\section{Algebraic Singularities and the Transfer Theorem}

Not every generating function arising from combinatorics is rational. The Catalan generating function $C(z) = (1 - \sqrt{1-4z})/(2z)$ has a square-root branch point, which is neither a pole nor a removable singularity. Binary trees, general planar maps, and most classes defined by functional equations of the form $T(z) = z \cdot \phi(T(z))$ give rise to algebraic singularities of this type \cite{Drmota2009}. The appropriate tool is the Flajolet--Odlyzko transfer theorem.

First we need a geometric notion. Intuitively, the challenge near a singularity $\rho$ is that $A$ may fail to be analytic exactly at $\rho$ while remaining analytic on a region extending slightly beyond the circle $|z| = \rho$ in all other directions. The $\Delta$-domain formalizes this.

\begin{definition}
For $\rho > 0$, $\phi \in (0, \pi/2)$, and $\eta > \rho$, the \emph{$\Delta$-domain at $\rho$} is the set
\[
\Delta(\rho, \phi, \eta) = \{ z \in \mathbb{C} : |z| < \eta,\; z \ne \rho,\; |\arg(z - \rho)| > \phi \}.
\]
We say $A$ is \emph{$\Delta$-analytic at $\rho$} if it is analytic in some $\Delta(\rho, \phi, \eta)$.
\end{definition}

Concretely, a $\Delta$-domain is a disk of radius $\eta > \rho$ centered at the origin, with a small wedge-shaped slit removed around the positive real ray at and beyond $\rho$. The function $A$ need not extend to $\rho$ itself, but it must be analytic in the slit disk. One may picture shining a narrow flashlight from the origin outward: the beam extends slightly past $\rho$ in all directions except along the slit, where $\rho$ itself blocks the view.

\begin{theorem}[Flajolet--Odlyzko Transfer Theorem, \cite{FlajoletOdlyzko1990}]
Let $\rho > 0$ and suppose $A(z)$ is $\Delta$-analytic at $\rho$. Suppose further that as $z \to \rho$ within the $\Delta$-domain,
\[
A(z) = C \left(1 - \frac{z}{\rho}\right)^{-\alpha} (1 + o(1))
\]
for some $\alpha \in \mathbb{R} \setminus \mathbb{Z}_{\le 0}$ and $C \ne 0$. Then
\[
[z^n] A(z) \sim \frac{C}{\Gamma(\alpha)}\, n^{\alpha - 1}\, \rho^{-n}
\]
as $n \to \infty$.
\end{theorem}

The proof is a beautiful application of contour deformation from Ch2, and we sketch it here to give the reader confidence that the formula is not magic.
Start with the Cauchy coefficient formula from Chapter~2:
\[
  a_n = \frac{1}{2\pi i} \oint_{|z|=r} \frac{A(z)}{z^{n+1}}\, dz,
\]
where the contour is a small circle of radius $r < \rho$ centered at the origin.
By the contour deformation principle (Chapter~2), we may deform this circle outward
to any path that stays within the region where $A$ is holomorphic.
Because $A$ is $\Delta$-analytic at $\rho$, we can push the contour past the circle
$|z| = \rho$ everywhere \emph{except} near the singularity, wrapping it tightly
around the branch cut or pole at $z = \rho$; this is the \emph{Hankel contour}.
The Hankel contour consists of a large outer arc plus two segments running along
opposite sides of the branch cut connecting $\rho$ to infinity.
By the ML bound from Chapter~2, the contribution from the large outer arc
is at most (length) $\times$ (maximum of $|A(z)/z^{n+1}|$ on the arc),
and since $|z^{n+1}|$ grows exponentially in $n$ on the arc, this contribution
decays exponentially and is negligible.
The two segments hugging the singularity are what survive, and on them
the integrand is governed entirely by the local behavior $A(z) \approx C(1-z/\rho)^{-\alpha}$.
Substituting $z = \rho(1 - t/n)$ and rescaling, the integral over the two segments
reduces to an Euler-type integral that evaluates to $\int_0^\infty t^{\alpha-1}e^{-t}\,dt = \Gamma(\alpha)$,
with the $\rho^{-n}$ and $n^{\alpha-1}$ factors emerging from the change of variables.
Dividing out $\Gamma(\alpha)$ produces the stated formula.
Full details appear in~\cite{FlajoletSedgewick2009}, Chapter~VI.

What matters for applications is the structural decomposition of the conclusion:

\begin{enumerate}
\item The exponential factor $\rho^{-n}$ is determined entirely by the \emph{location} of the singularity, not its type. Moving the singularity closer to or farther from the origin changes the exponential growth rate.
\item The polynomial factor $n^{\alpha - 1}$ is determined by the \emph{order of contact} of $A$ with its singularity. A large positive $\alpha$ means $A$ blows up strongly at $\rho$, and the coefficients grow correspondingly fast. A negative $\alpha$ means $A$ vanishes at $\rho$, and the coefficients decay polynomially relative to $\rho^{-n}$.
\item The constant $C/\Gamma(\alpha)$ is fixed by the local behavior of $A$ near $\rho$; it is independent of $n$ and encodes global information about the combinatorial class through the single number $C$.
\end{enumerate}

The $\Delta$-analyticity hypothesis is essential. It rules out situations where a second dominant singularity (not at $\rho$) interferes. In the aperiodic setting, one typically verifies $\Delta$-analyticity by confirming that all other singularities of $A$ lie strictly outside $|z| = \rho$ or can be analytically continued around.

\section{Catalogue of Exponents}

The transfer theorem applies across a range of singularity types. We record the most important cases.

\paragraph{Poles.} If $\alpha = k$ is a positive integer, then $(1 - z/\rho)^{-k}$ has a pole of order $k$ at $\rho$. The gamma function value is $\Gamma(k) = (k-1)!$, so
\[
[z^n] \frac{C}{(1 - z/\rho)^k} \sim \frac{C}{(k-1)!}\, n^{k-1}\, \rho^{-n}.
\]
For $k = 1$ (simple pole) this gives $a_n \sim C \rho^{-n}$ with no polynomial factor. For $k = 2$ (double pole), $a_n \sim Cn\rho^{-n}$. This is consistent with the partial fraction analysis of Section~3.

\paragraph{Square-root singularity.} If $A(z) \sim C\sqrt{1 - z/\rho}$ near $\rho$, then $\alpha = -1/2$ and $A$ vanishes like a square root. Since $\Gamma(-1/2) = -2\sqrt{\pi}$,
\[
a_n \sim \frac{C}{-2\sqrt{\pi}}\, n^{-3/2}\, \rho^{-n} = -\frac{C}{2\sqrt{\pi}}\, n^{-3/2}\, \rho^{-n}.
\]
The $n^{-3/2}$ decay relative to $\rho^{-n}$ is characteristic; the negative sign in $\Gamma(-1/2)$ combines with the sign of $C$ to give a positive coefficient sequence when $C < 0$ (which is the typical case when $A$ approaches $0$ at $\rho$ from below).

\paragraph{Inverse square-root singularity.} If $A(z) \sim C/\sqrt{1 - z/\rho}$, then $\alpha = 1/2$. Since $\Gamma(1/2) = \sqrt{\pi}$,
\[
a_n \sim \frac{C}{\sqrt{\pi}}\, n^{-1/2}\, \rho^{-n}.
\]
The $n^{-1/2}$ polynomial prefactor is the signature of an inverse square-root singularity. This arises, for instance, in counting lattice paths with unconstrained endpoint.

\paragraph{Logarithmic singularity.} The function $\log(1/(1 - z/\rho))$ is a limiting case outside the strict hypotheses of the theorem as stated (since $\alpha \to 0$ as $z \to \rho$). A direct computation gives $[z^n] \log(1/(1-z/\rho)) = \rho^{-n}/n$, a $1/n$ decay that sits between the $\alpha \to 0$ and $\alpha \to 1$ cases.

\section{Worked Example: Catalan Numbers}

\begin{example}[Catalan numbers]
In Chapter~1 we derived the generating function for the Catalan numbers $C_n = \binom{2n}{n}/(n+1)$:
\[
C(z) = \frac{1 - \sqrt{1-4z}}{2z}.
\]
The function $\sqrt{1-4z}$ has a branch point at $z = 1/4$, which is the dominant singularity of $C(z)$. All other singularities (the apparent pole at $z=0$ is removable) lie no closer to the origin. Thus $\rho = 1/4$.

To apply the transfer theorem, expand $C(z)$ near $z = 1/4$. Let $u = 1 - 4z$, so $u \to 0$ as $z \to 1/4$. Then
\[
C(z) = \frac{1 - \sqrt{u}}{2z} = \frac{1 - \sqrt{u}}{(1 - u)/2} = \frac{2(1 - \sqrt{u})}{1 - u}.
\]
Expanding $(1-u)^{-1} = 1 + u + O(u^2)$ and keeping terms through $O(\sqrt{u})$:
\[
C(z) = 2(1 - \sqrt{u})(1 + u + O(u^2)) = 2 - 2\sqrt{u} + O(u) = 2 - 2\sqrt{1-4z} + O(1 - 4z).
\]
Since $\sqrt{1-4z} = \sqrt{4(1/4 - z)} = 2\sqrt{1/4 - z}$, we can also write $\sqrt{1-4z} = \sqrt{4}\sqrt{1 - z/\rho} = 2(1-z/\rho)^{1/2}$, so
\[
C(z) = 2 - 4(1 - z/\rho)^{1/2} + O(1 - z/\rho).
\]
The constant term $2 = C(1/4)$ does not contribute to the asymptotics of coefficients. The singular part is $-4(1 - z/\rho)^{1/2}$, corresponding to $\alpha = 1/2$ and $C = -4$. However, this is not the form $(1-z/\rho)^{-\alpha}$; rather it is $(1-z/\rho)^{+1/2}$, corresponding to $-\alpha = 1/2$, i.e., $\alpha = -1/2$ in the theorem's notation. We have $C = -4$ and $\alpha = -1/2$.

A naive first attempt at applying the transfer theorem might use $C_{\mathrm{sing}} = -4$, but this would give $C_n \sim 2 \cdot 4^n/(\sqrt{\pi}\, n^{3/2})$, which is off by a factor of $2$ from the classical asymptotic. The error lies in the change of variables. The full expansion is $C(z) = 2 - 2\sqrt{1-4z} + O(1-4z)$, and the coefficient of $(1-4z)^{1/2}$ is $-2$, not $-4$. Since $1-4z = 4(1/4-z) = 4\rho(1 - z/\rho)$, we have $\sqrt{1-4z} = 2\sqrt{\rho}\sqrt{1-z/\rho} = \sqrt{1-z/\rho}$ (because $\rho = 1/4$, so $\sqrt{4\rho} = 1$). Therefore $C(z) \sim 2 - 2(1-z/\rho)^{1/2}$, giving $C_{\mathrm{sing}} = -2$ and $\alpha = -1/2$. Applying the transfer theorem:
\[
C_n \sim \frac{-2}{\Gamma(-1/2)}\, n^{-3/2}\, 4^n = \frac{-2}{-2\sqrt{\pi}}\, n^{-3/2}\, 4^n = \frac{4^n}{\sqrt{\pi}\, n^{3/2}}.
\]
This agrees exactly with the asymptotic derived from Stirling's formula applied to the exact formula $C_n = \binom{2n}{n}/(n+1)$. The moral is that the change of variables from $(1-4z)$ to $(1-z/\rho)$ must be tracked carefully; a factor of $\sqrt{4\rho}$ is easy to drop. The $n^{-3/2} \cdot 4^n$ form is the prototype of what Chapter~5 will call the \emph{$n^{-3/2}$ law}.
\end{example}

\section{Summary and Preview}

The Flajolet--Odlyzko transfer theorem completes the central algorithmic workflow of analytic combinatorics. Given a combinatorial class $\mathcal{A}$: (i) apply the symbolic method (Chapter~3) to translate the structural decomposition of $\mathcal{A}$ into a functional equation for the generating function $A(z)$; (ii) analyze that equation to locate the dominant singularity $\rho$ and determine the local exponent $\alpha$; (iii) read off the asymptotic $a_n \sim (C/\Gamma(\alpha))\, n^{\alpha-1}\, \rho^{-n}$ directly from the singularity data. No other computation is required.

The power of this approach lies in what it unifies. Rational GFs yield polynomial-times-exponential asymptotics via poles; algebraic GFs yield the same form with non-integer exponents via branch points. Both cases are handled by the single theorem~\cite{FlajoletOdlyzko1990}, with the difference in exponent type $\alpha$ being the only distinction.

Chapter~5 will examine the class of context-free grammars in detail. A fundamental theorem of analytic combinatorics states that any unambiguous context-free grammar gives rise to an algebraic generating function, and that in the aperiodic case the dominant singularity is always a square-root branch point, hence $\alpha = -1/2$ universally. This means the coefficient of any unambiguous context-free language satisfies $a_n \sim C \cdot n^{-3/2} \cdot \rho^{-n}$ (with $\rho$ and $C$ depending on the grammar, but the exponent $-3/2$ being universal), a remarkable structural prediction. The same $n^{-3/2}$ law turns out to be intimately connected with the length distribution of sentences generated by language models, a connection we will develop in Parts~III--V.

\section*{Further Reading}

\begin{itemize}
  \item \textbf{P. Flajolet and A.\,M. Odlyzko}, ``Singularity analysis of generating functions,'' \emph{SIAM J.\ Discrete Math.} (1990) \cite{FlajoletOdlyzko1990} --- The original paper proving the transfer theorem in the form used throughout this chapter; contains the Hankel-contour proof and the full catalogue of singular types, and is more concise than the book-length treatment in \cite{FlajoletSedgewick2009}.

  \item \textbf{P. Flajolet and R. Sedgewick}, \emph{Analytic Combinatorics} (Cambridge University Press, 2009) \cite{FlajoletSedgewick2009}, Chapter~VI --- The textbook development of singularity analysis, including the $\Delta$-domain framework, a detailed treatment of the Hankel contour argument, and worked examples covering random trees, urns, and maps.

  \item \textbf{M. Drmota}, \emph{Random Trees} (Springer, Vienna, 2009) \cite{Drmota2009} --- A monograph on random trees whose first half is an accessible introduction to the singularity analysis of algebraic and D-finite generating functions; Chapter~2 is especially useful for understanding how the characteristic system of the present chapter extends to systems of equations.

  \item \textbf{A. Odlyzko}, ``Asymptotic enumeration methods,'' in \emph{Handbook of Combinatorics}, Vol.~2 (Elsevier, 1995) --- A comprehensive survey of coefficient-extraction techniques from generating functions, covering the transfer theorem, saddle-point methods, and the Darboux method; excellent for seeing the transfer theorem in context alongside other asymptotic tools.

  \item \textbf{C. Banderier and M. Drmota}, ``Formulae and asymptotics for coefficients of algebraic functions,'' \emph{Combinatorics, Probability and Computing} (2015) \cite{BanderierDrmota2015} --- Proves universal limit laws for coefficients of algebraic generating functions beyond the leading term; relevant for understanding why the $n^{-3/2}$ exponent governs not just asymptotics but also distributional limits of combinatorial parameters.
\end{itemize}
